%%%%%%%%%%  scan idx 320  |  folio 307  %%%%%%%%%%
est un
isomorphisme.

Pour le voir qu'il est injectif, on utilise la suite d'homomorphismes canoniques~:
\[
  \Hom(\Ghat,H\circ g)\longrightarrow
  \Hom(\Ghat\circ i,H)\longrightarrow
  \Hom(\Gg,H\circ g)\longrightarrow
  \Hom(\Ghat,H\circ g),
\]
la dernière flèche venant de ce que \(\Ghat\) est un sous-foncteur de \(\Gg\). Pour vérifier que l'homomorphisme composé est l'identité, on remarque qu'un morphisme de foncteurs~: \(\Ghat\longrightarrow H\circ g\) est déterminé par sa valeur pour les objets de \(\mathcal C\) dont la filtration est grossière~; en effet, soit \(A\in\operatorname{Ob}(\mathcal C)\), et \(A_g\) l'algèbre \(A\) munie de la filtration grossière~; le morphisme canonique \(A\to A_g\) définit un diagramme commutatif~:
\[
\begin{tikzcd}
\Ghat(A) \arrow[r,"\Phi_A"] \arrow[d] &
H\circ g(A) \arrow[d]\\
\Ghat(A_g) \arrow[r,"\Phi_{A_g}"] &
H\circ g(A_g),
\end{tikzcd}
\]
montrant que \(\Phi_A\) est la restriction de \(\Phi_{A_g}\) à \(\Ghat(A)\).

On considère ensuite l'homomorphisme composé~:
\[
  \Hom(\Ghat\circ i,H)\longrightarrow
  \Hom(\Gg,H\circ g)\longrightarrow
  \Hom(\Ghat,H\circ g)\longrightarrow
  \Hom(\Ghat\circ i,H),
\]
et il est clair que c'est l'identité, d'où le résultat annoncé.

\textbf{Théorème 1.10.} Soit \(H\) un foncteur sur la catégorie \(\mathcal C_0\) des \(K\)-algèbres, à valeurs dans la catégorie des groupes abéliens filtrés séparés complets, et commutant aux invariants sous un groupe fini. Alors l'application de \(\Hom_{\Abfil}(\Ghat\circ i,H)\) dans \(H(K[X])\) qui, à un morphisme de foncteurs, associe
